% abstract.tex — ESWA abstract (under 250 words)

Online hiring platforms rank millions of resume--job pairs, yet return a single opaque score with no breakdown of which fields drove the rank, so recruiters cannot justify shortlists and candidates cannot act on what would change theirs. We present \JobMatch, an auditable, calibrated, and explainable recommendation methodology for recruitment (implemented as a multi-agent system) that combines four contributions. First, a composite score that decomposes into six documented channels (semantic similarity, skill overlap, title fit, experience tier, compensation fit, remote policy), giving component-level explanations without post-hoc attribution; alongside it, a graded relation-aware skill matcher that, on a de-circularized benchmark, collapses orthographic/synonym variants to exact (recall 1.00) yet keeps most non-equivalent look-alikes distinct (seven of eight, with one flagged catalog-data merge). Second, a calibrated-confidence layer reported with an honest calibration--discrimination trade-off: Platt scaling attains a low equal-width ECE (0.019) but near-base-rate discrimination (adaptive ECE 0.084, AUC 0.76), whereas beta calibration preserves discrimination (adaptive ECE 0.009, AUC 0.96); we keep Platt as the frozen default and recommend beta. Third, structural \emph{and} mechanistic explanation-faithfulness checks (the latter: does removing the top-attributed channel displace the top-ranked item), plus a 50-pair counterfactual probe. Fourth, a reproducible, regression-gated open-source artifact (frozen 30-resume/15-job/47-label corpus, one-command evaluation).
Evaluation is deliberately controlled but small (30 resumes, 15 jobs, 47 human-labeled pairs, single annotator, positive-only labels), so results establish methodological \emph{feasibility}, not population-level ranking superiority: the composite reaches nDCG@5 = 0.949 but its gain over the semantic baseline (0.878) is not significant ($p = 0.10$; nothing survives Holm correction at $n = 30$), and a 25-configuration protocol-gated search finds no better ranker. Decomposing the skill channel, an asymmetric coverage form drives the synthetic gain (Jaccard $0.917 \rightarrow 0.949$) while the relation-aware credit --- the novel component --- helps only on the real corpus (six of thirty queries improve, none worsen; sign-test $p = 0.03$), which we report as a directional signal rather than a broad effect.

